%---------------------------------------------------------------------------%
%->> Frontmatter
%---------------------------------------------------------------------------%
%-
%-> 生成封面
%-
\maketitle% 生成中文封面
{\hfuzz=10pt\MAKETITLE}% 生成英文封面
%-
%-> 作者声明
%-
\makedeclaration% 生成声明页
%-
%-> 中文摘要
%-
\intobmk\chapter*{摘要}% 显示在书签但不显示在目录
\setcounter{page}{1}% 开始页码
\pagenumbering{Roman}% 页码符号

大语言模型在数学推理等可验证任务上的后训练，越来越依赖基于组采样的强化学习方法。然而，固定小组采样会导致一个经常被忽视的统计问题：当同一 prompt 采样得到的回答全对或全错时，组内奖励方差退化，进而造成零梯度更新。本文将这一现象概括为\textit{signal collapse}，并指出其根源并不完全在于题目是否可学，而在于训练预算与 prompt 难度之间存在显著错配。

围绕这一问题，本文基于 \texttt{Qwen2.5-Math-1.5B base} 在 \texttt{OpenR1-Math-220k} 上的采样分布与训练过程切片数据，系统分析了固定预算、自适应固定块预算和几何块预算的差异。研究表明：训练数据内部存在显著的难度长尾结构；即使组大小提升到 $32$，训练前后仍分别有 35.36\% 与 26.41\% 的 prompt 发生 signal collapse；在此背景下，自适应采样较固定预算具有更低的信号缺失率，而几何块自适应策略在保持统计效果的同时，更适合后续与 KV Cache 复用等系统优化结合。进一步地，本文比较了“有对有错就停”与“只要有对就停”两类退出规则，发现正样本优先规则在现有预算下具有更高的预算效率。

本文的工作说明，在 reasoning 强化学习后训练中，采样预算本身就是训练信号设计的一部分。与其仅从损失函数角度理解训练稳定性，不如将 prompt 级预算分配、退出规则和系统批次组织作为统一问题加以设计。

\keywords{大语言模型，强化学习后训练，自适应采样，信号坍塌，数学推理}% 中文关键词
%-
%-> 英文摘要
%-
\intobmk\chapter*{Abstract}% 显示在书签但不显示在目录

Reinforcement learning post-training for large language models has recently become a major approach for improving mathematical reasoning. A practical difficulty of group-based optimization, however, is that fixed small-group sampling often yields all-correct or all-incorrect response sets for the same prompt, causing the within-group reward variance to vanish and the gradient signal to collapse. This thesis studies the problem from the perspective of prompt-level sampling allocation and argues that many zero-signal cases are induced by budget mismatch rather than the intrinsic unlearnability of prompts.

Using prompt-level repeated sampling statistics and checkpoint slices from the Qwen2.5-Math-1.5B base model trained on OpenR1-Math-220k, this thesis analyzes the distribution of prompt difficulty, the evolution of signal collapse during training, and the trade-off among fixed-budget, fixed-block adaptive, and geometric-block adaptive sampling strategies. The results show that the sampled training data exhibits a pronounced long-tail difficulty structure; signal collapse remains substantial even with group size $32$; adaptive sampling consistently outperforms comparable fixed budgets; and geometric-block scheduling achieves nearly the same statistical effectiveness as fixed-block scheduling while offering a more system-friendly batch structure. Moreover, under the current budget, a positive-first stopping rule is more cost-effective than a balanced stopping rule.

These findings suggest that sampling allocation should be viewed as an integral part of signal design in reasoning-oriented reinforcement learning, together with the stopping rule and systems-level batch organization.

\KEYWORDS{large language models, reinforcement learning post-training, adaptive sampling, signal collapse, mathematical reasoning}% 英文关键词
%---------------------------------------------------------------------------%
